% ██████████████████████████████████████████████████████████████████
%  CHAPTER 3 — METHODOLOGY
% ██████████████████████████████████████████████████████████████████
\chapter{Methodology}
\label{ch:methodology}

\section{Research Strategy}
\label{sec:methodology_strategy}

This thesis adopts a \emph{constructive--empirical} methodology that combines a controlled experimental study (to characterise the XAI degradation profile) with a design-science contribution (to propose and implement the Maturity-Aware AI Routing architecture). The two strands are mutually reinforcing: the experimental findings provide the empirical justification for the architectural threshold, and the architectural implementation provides the operational vehicle through which the findings are translated into a deployable artefact.

\section{Datasets}
\label{sec:methodology_data}

The Late-Payment exemplar is built on the publicly available Brazilian Olist e-commerce corpus, which provides $96{,}460$ joined order--payment--item records after the standard preprocessing pipeline (orders $\bowtie$ payments $\bowtie$ items). The binary target is constructed by treating an order whose actual delivery date exceeds the estimated delivery date as a positive (\emph{late}) instance, yielding a base positive rate of $8.1\%$. Fourteen features are engineered, spanning monetary amounts, item counts, payment-type one-hots, temporal lags (approval delay, carrier delay), and purchase-time signals.

\section{Sparsity Simulation Protocol}
\label{sec:methodology_sparsity}

Five tenant slices of geometrically decreasing volume are generated through stratified random sampling without replacement:

\begin{itemize}
  \item Tenant A (\emph{data-rich}) --- $n = 50{,}000$
  \item Tenant B (\emph{medium})    --- $n = 5{,}000$
  \item Tenant C (\emph{sparse})    --- $n = 1{,}000$
  \item Tenant D (\emph{very sparse}) --- $n = 200$
  \item Tenant E (\emph{cold-start}) --- $n = 50$
\end{itemize}

Stratification preserves the global positive-class proportion within each slice. All slices are drawn with a fixed random seed ($\mathtt{SEED} = 42$) for reproducibility.

\section{Model Training and Evaluation}
\label{sec:methodology_training}

Each tenant slice is split $80/20$ into train/test partitions, stratified on the late-payment label. A \texttt{GradientBoostingClassifier} (scikit-learn defaults, $\mathtt{random\_state}=42$) is trained on the train partition; predictive performance is reported on the held-out test partition using AUC-ROC and the macro-averaged F1-score.

\section{Explanation Generation}
\label{sec:methodology_explanation}

Global feature importances are computed via \texttt{shap.TreeExplainer} applied to the trained classifier and evaluated on the test partition. Local explanations are generated with \texttt{lime.lime\_tabular.LimeTabularExplainer} using a representative background drawn from the train partition. Per-instance LIME stability is assessed by drawing five representative instances and repeating each explanation five times with distinct random seeds.

\section{Metric Definitions}
\label{sec:methodology_metrics}

\paragraph{Spearman Rank Fidelity ($\rho$).} For each tenant, the global SHAP importance vector is ranked and compared against Tenant A's ranking via Spearman's $\rho$. Values close to $1.0$ indicate strong agreement with the data-rich reference; values approaching $0$ indicate explanation drift.

\paragraph{LIME Top-$k$ Jaccard Stability.} For each of the five representative instances, the top-$5$ features from each of the five repeats are extracted, and the average pairwise Jaccard overlap is reported.

\paragraph{SHAP Shannon Entropy ($H$).} The normalised mean absolute SHAP importances are treated as a probability distribution over features; lower entropy indicates concentration on a small number of features (a symptom of grasping-at-straws behaviour in sparse regimes).

\paragraph{SHAP Absolute-Value Variance.} The variance of $\lvert \mathrm{SHAP}_{ij}\rvert$ across all (instance, feature) pairs in the test partition; higher variance indicates more volatile per-instance attributions.

\section{Architectural Construction}
\label{sec:methodology_arch}

The Maturity-Aware AI Routing pattern is implemented as a new method, \texttt{predict\_or\_heuristic}, on the existing \texttt{LatePaymentPredictor} service class in \texttt{erp/ai/services/finance.py}. The method counts the tenant's closed-status invoices, compares this count against the class-level constant \texttt{COLD\_START\_INVOICE\_THRESHOLD = 1000}, and returns either a local probability (with mode \texttt{"local\_prediction"}) or the platform-wide heuristic (with mode \texttt{"global\_heuristic"}) loaded from the versioned JSON artefact \texttt{thesis/results/sparsity\_experiment/cold\_start\_heuristic.json}. The existing \texttt{predict()} signature is preserved for backward compatibility.

\section{Reproducibility}
\label{sec:methodology_reproducibility}

All experimental code, configuration, and outputs are version-controlled under \texttt{thesis/scripts/} and \texttt{thesis/results/sparsity\_experiment/}. The end-to-end pipeline can be re-executed with a single command using the project's pinned Python virtual environment.
